\scanpage{59}
\mspage{29}
prète comme le groupe \inserted{\unsure{topologique|topologisons}} des $U$-autom du schéma $\widetilde U$, et $E_U\longrightarrow E_K$ comme l'hom de passage au quotient induit sur $\bar K$, identifié à la clôture alg. de $K$ dans $H^0(\widetilde U,\mathcal O_{\widetilde U})$, donc $E_U$ opère sur $\operatorname{Spec}\bar K$ de façon naturelle. Une section de $E_U$ sur $E_K$ est donc une action de $E_K$ sur $\widetilde U$, compatible avec son action sur $\unsure{\widetilde U|\bar U}$, et continue dans un sens convenable (sans doute la conv. \unsure{simple-discrète} sur les \deleted{$\unsure{\widetilde U_i}$} $\bar U_i$ finis sur $\bar U$ entre $\widetilde U$ et $\bar U$\ldots).

Considérons alors la

\uline{Conjecture B} Toute telle action de $\Gamma$ sur $\widetilde U$ admet \uline{dans $X'$}$=\widehat{\widetilde U}$ un pt fixe et un seul.
\margin{\figure{\begin{tikzpicture}\draw (0,0)--(0,-3.1);\end{tikzpicture}}}
\margin{et $\unsure{=|\overline m}$ l'action induite de $\Gamma^\circ$ doit avoir un pt fixe et un seul --- cf plus bas.}

Ceci signifie donc

a) S'il y a un pt fixe à distance finie i.e. $\widetilde x\in\widetilde U^\Gamma$, alors 1°) L'image de $\widetilde x$ dans $U$ est uniquement déterminée --- c'est essentiellement le \unsure{th.3} du $\S1$ (deux pts distincts de $U(K)$ définissent des classes de conj. de scindages distinctes) et
\margin{58}
